\subsection*{3.3 微分中值定理}

本节的主要方法是构造辅助函数，使它在若干点取相同的值，再反复应用 Rolle 中值定理。对较复杂的表达式，可先减去满足边界条件的简单函数。

\ex{3.3.1}设 $f$ 在 $[0,\pi]$ 上二阶可导，且 $|f(x)|\leq\min\{1-\cos x,\sin x\}$。证明存在 $\xi\in(0,\pi)$，使 $f''(\xi)=f(\xi)$。
\pf 由条件得 $f(0)=f(\pi)=0$，并由 $|f(x)|/x\leq(1-\cos x)/x\to0$ 得 $f'(0)=0$。令 $g(x)=e^{-x}f(x)$，则 $g(0)=g(\pi)=0$，故存在 $\xi_1\in(0,\pi)$ 使 $g'(\xi_1)=0$。再令 $h(x)=e^xg'(x)=f'(x)-f(x)$，有 $h(0)=h(\xi_1)=0$。又由 Rolle 定理，存在 $\xi\in(0,\xi_1)$ 使
\[
h'(\xi)=f''(\xi)-f(\xi)=0.\qquad\square
\]
\noindent\textit{校勘：原书将结论写为 $\xi\in(0,1)$，但原证明只能得到 $\xi\in(0,\pi)$；这里按可证明的结论排版。}

\ex{3.3.2}设 $f$ 在 $[1,+\infty)$ 上连续可微，$f(1)=1$，且
\[
f'(x)=\frac1{x^2+f(x)^2}\quad(x\geq1).
\]
估计 $A=\sup_{x\geq1}f(x)$。
\sol 因 $f'>0$，$f$ 严格递增；又因 $f(x)>1$（$x>1$），
\[
0<f(x)-1=\int_1^x\frac{\mathrm dt}{t^2+f(t)^2}
<\int_1^\infty\frac{\mathrm dt}{t^2+1}=\frac\pi4.
\]
因此 $A=f(+\infty)<1+\pi/4$。再由 $f(x)<A$ 得
\[
A-1>\int_1^\infty\frac{\mathrm dx}{x^2+A^2}
=\frac1A\left(\frac\pi2-\arctan\frac1A\right).
\]
特别地，利用 $f^2f'>1/(x^2+1)$，有
\[
\frac{A^3-1}{3}=\int_1^\infty f^2f'\,\mathrm dx
>\frac\pi4,
\qquad
\sqrt[3]{1+\frac{3\pi}{4}}<A<1+\frac\pi4.
\]
设 $q(t)=t^2-t+\arctan(1/t)-\pi/2$，则上述隐式下界还给出 $A>A_0$，其中 $A_0>1$ 是 $q(A_0)=0$ 的唯一根，约为 $1.62$。

\ex{3.3.3}设 $f$ 在 $[-3,1]$ 上二阶可导，且 $f(-3)=-2$、$f(0)=1$、$f(1)=10/3$。证明存在 $\xi\in(-3,1)$ 使 $f''(\xi)=2\xi+2$。
\pf 取 $p(x)=x^2+x^3/3+x+1$，则 $p(-3)=-2$、$p(0)=1$、$p(1)=10/3$，且 $p''(x)=2x+2$。令 $F=f-p$，由 $F(-3)=F(0)=F(1)=0$，连续两次应用 Rolle 定理，得某个 $\xi\in(-3,1)$ 满足 $F''(\xi)=0$。\hfill$\square$

\ex{3.3.4}设 $f$ 在 $[-1,2]$ 上二阶可导，且 $f(-1)=2-e^{-1}$、$f(0)=1$、$f(2)=2e^2-1$。证明存在 $\xi\in(-1,2)$ 使
\[
f''(\xi)=f(\xi)+2e^\xi+\xi-1.
\]
\pf 取 $p(x)=xe^x-x+1$，则 $p''-p=2e^x+x-1$，并且 $p$ 在 $-1,0,2$ 处与 $f$ 取值相同。令 $F=f-p$。函数 $e^{-x}F(x)$ 在三处为零，故它的导数在两处为零，即 $G(x)=F'(x)-F(x)$ 在两处为零。再对 $e^xG(x)$ 应用 Rolle 定理，得某个 $\xi\in(-1,2)$ 满足
\[
0=(e^xG(x))'\big|_{x=\xi}=e^\xi\bigl(F''(\xi)-F(\xi)\bigr),
\]
结论随之成立。\hfill$\square$

\ex{3.3.5}设 $f$ 在 $[a,b]$ 上二阶可导，且 $f'(a)=f'(b)=0$。证明：
\begin{enumerate}[label=（\roman*）,leftmargin=3em]
\item 存在 $\xi\in(a,b)$ 使 $|f''(\xi)|\geq4|f(b)-f(a)|/(b-a)^2$；
\item 若 $f(b)\ne f(a)$，可将上式中的 $\geq$ 改为 $>$。
\end{enumerate}
\pf 记 $L=b-a$，$D=|f(b)-f(a)|$。若 $D=0$，第一条显然。若 $D>0$，设 $K=4D/L^2$。反设所有内点均有 $|f''|\leq K$，由中值定理及两个端点的导数条件，
\[
|f'(x)|\leq K\min\{x-a,b-x\}\quad(a<x<b).
\]
于是
\[
D=\left|\int_a^bf'(x)\,\mathrm dx\right|
\leq K\int_a^b\min\{x-a,b-x\}\,\mathrm dx=\frac{KL^2}{4}=D.
\]
若处处取等号，$f'$ 必等于同号的 $K\min\{x-a,b-x\}$；但此函数在 $x=(a+b)/2$ 不可导，与 $f$ 二阶可导矛盾。故某处 $|f''|>K$，两条结论均得证。\hfill$\square$

\ex{3.3.7}设 $f$ 在 $[a,b]$ 上 $(n+1)$ 阶可导，$f^{(k)}(a)=f^{(k)}(b)=0$（$1\leq k\leq n$），且 $f(b)\ne f(a)$。证明存在 $\xi\in(a,b)$ 使
\[
|f^{(n+1)}(\xi)|>
\frac{2^n(n+1)!}{(b-a)^{n+1}}|f(b)-f(a)|.
\]

\ex{3.3.8}设 $f$ 二阶可导，且 $f(0)=0$、$f'(0)=1$、$f(\pi/4)=1$。证明存在 $\xi\in(0,\pi/4)$ 使 $f''(\xi)=2f(\xi)f'(\xi)$。
\pf 函数 $H=f''-2ff'$ 是导函数 $(f'-f^2)'$，故具有介值性。若它在 $(0,\pi/4)$ 上无零点，只可能恒正或恒负。若恒正，则 $f'-f^2>1$，从而
\[
(\arctan f)'=\frac{f'}{1+f^2}>1.
\]
在 $[0,\pi/4]$ 上积分得到
\[
\arctan f(\pi/4)>\frac\pi4,
\]
矛盾。恒负时得到反向的严格不等式，也与 $\arctan f(\pi/4)=\pi/4$ 矛盾。\hfill$\square$

\ex{3.3.9}设 $f$ 在 $\mathbb R$ 上三阶可导，且 $f,f',f'',f'''>0$、$f'''\leq f$。证明 $f'(x)<2f(x)$ 对所有 $x$ 成立。
\pf 固定 $x$。对 $t<0$ 用 Taylor 公式及 $f(x+t)>0$，得
\[
0<f(x+t)<f(x)+f'(x)t+\tfrac12f''(x)t^2.
\]
此二次式对 $t\geq0$ 也为正，故判别式为负，即 $(f')^2<2ff''$。同样对 $f'(x+t)>0$ 展开，并由 $f'''(\eta)\leq f(\eta)<f(x)$（$t<0$），可得 $(f'')^2<2ff'$。将两式合用，
\[
(f')^4<4f^2(f'')^2<8f^3f',
\]
由于 $f'>0$，故 $(f')^3<8f^3$，即 $f'<2f$。\hfill$\square$

\thm{3.3.1}设 $f$ 在 $[x_1,x_k]$ 上 $n$ 次可导，并在 $x_1<\cdots<x_k$ 处分别具有连续排列的 $n_1,\ldots,n_k$ 个零插值条件，且 $n_1+\cdots+n_k=n+1$。若实系数多项式
\[
\lambda^n+a_{n-1}\lambda^{n-1}+\cdots+a_1\lambda+a_0
\]
的根全为实数，则存在 $\xi\in(x_1,x_k)$ 使
\[
f^{(n)}(\xi)+a_{n-1}f^{(n-1)}(\xi)+\cdots+a_1f'(\xi)+a_0f(\xi)=0.
\]
这里零插值条件指 $f(x_i)=f'(x_i)=\cdots=f^{(n_i-1)}(x_i)=0$，并须保证端点所需的单侧导数存在。其思路是把左边依次写成 $(D-\lambda_1)\cdots(D-\lambda_n)f$，每一步乘指数因子后应用 Rolle 定理。

\ex{3.3.10}设 $f$ 在 $[-1,1]$ 上三阶可导，$f(-1)=0$、$f(1)=1$、$f'(0)=0$。证明存在 $\xi\in(-1,1)$ 使 $f'''(\xi)=3$。
\pf 令
\[
p(x)=\tfrac12x^3+\bigl(\tfrac12-f(0)\bigr)x^2+f(0),\qquad F=f-p.
\]
则 $F(-1)=F(0)=F(1)=F'(0)=0$。由 Rolle 定理，$F'$ 在 $(-1,0)$、$(0,1)$ 各有一零点，加上 $F'(0)=0$；继续两次应用 Rolle 定理，得某个 $\xi$ 使 $F'''(\xi)=0$。因 $p'''=3$，结论成立。\hfill$\square$

\thm{3.3.2}设 $f$ 在 $(a,b)$ 内二阶可导且有三个不同的零点。若 $t^2+pt+q$ 的两根为 $\alpha\pm i\beta$，其中 $\beta>0$，并且 $\beta(b-a)\leq\pi$，则存在 $\xi\in(a,b)$ 使
\[
f''(\xi)+pf'(\xi)+qf(\xi)=0.
\]
\pf 取 $c=(a+b)/2$，由 $\beta(b-a)\leq\pi$ 知 $\cos\beta(x-c)>0$（$a<x<b$）。函数
\[
H_1(x)=\frac{e^{-\alpha x}f(x)}{\cos\beta(x-c)}
\]
有三个不同零点，故其导数至少有两个零点。令
\[
H_2(x)=H_1'(x)\cos^2\beta(x-c)
=e^{-\alpha x}\bigl[(f'(x)-\alpha f(x))\cos\beta(x-c)+\beta f(x)\sin\beta(x-c)\bigr].
\]
$H_2$ 也有两个零点；再次应用 Rolle 定理，某处 $H_2'=0$。直接求导得
\[
H_2'(x)=e^{-\alpha x}\cos\beta(x-c)
\bigl[f''(x)-2\alpha f'(x)+(\alpha^2+\beta^2)f(x)\bigr].
\]
因 $p=-2\alpha$、$q=\alpha^2+\beta^2$，结论成立。\hfill$\square$

\ex{3.3.12}设 $f$ 在 $[-1,1]$ 上二阶可导，$f(-1)=f(0)=f(1)=0$。证明：
\begin{enumerate}[label=（\roman*）,leftmargin=3em]
\item 在 $(-1,0)$、$(0,1)$ 内各有一点 $\xi_i$ 满足 $f(\xi_i)=(1+\xi_i^2)f'(\xi_i)$；
\item 存在 $\eta\in(-1,1)$ 使 $(1-2\eta)f(\eta)=(1+\eta^2)^2f''(\eta)$。
\end{enumerate}
\pf 令 $g(x)=e^{-\arctan x}f(x)$。由于 $g(-1)=g(0)=g(1)=0$，Rolle 定理给出两点 $\xi_1\in(-1,0)$、$\xi_2\in(0,1)$ 满足 $g'(\xi_i)=0$，即第一条。再令
\[
H(x)=e^{\arctan x}\left(f'(x)-\frac{f(x)}{1+x^2}\right).
\]
因 $H(\xi_1)=H(\xi_2)=0$，存在 $\eta\in(\xi_1,\xi_2)$ 使 $H'(\eta)=0$。直接求导得
\[
H'(x)=e^{\arctan x}\left[f''(x)+\frac{(2x-1)f(x)}{(1+x^2)^2}\right],
\]
故第二条成立。\hfill$\square$
